\section{Summary}\label{summary}

We present eight R packages for cointegration testing and Granger
causality analysis: \texttt{cointsmall}, \texttt{xtcadfcoint},
\texttt{xtbreakcoint}, \texttt{hatemicoint}, \texttt{xtdhcoint},
\texttt{fcoint}, \texttt{makicoint}, and \texttt{caustests}. These
packages implement a wide range of cointegration tests---from
small-sample methods with as few as twelve observations to large panel
tests accommodating cross-sectional dependence, structural breaks, and
smooth Fourier approximations---as well as comprehensive Granger
causality tests including Fourier, quantile, and bootstrap variants. All
packages are open-source under the GPL-3 license.

\section{Statement of Need}\label{statement-of-need}

Cointegration analysis is fundamental to empirical economics,
underpinning the estimation of long-run equilibrium relationships
between non-stationary variables. While the \texttt{urca} package (Pfaff
2008) provides classical Engle-Granger and Johansen tests, and the
\texttt{cointReg} package offers cointegrating regression estimators,
the econometric literature has advanced substantially in several
directions that remain underserved in R:

\begin{enumerate}
\def\labelenumi{\arabic{enumi}.}
\tightlist
\item
  \textbf{Small-sample cointegration}: Standard tests perform poorly
  with fewer than 50 observations. The Trinh (2022) method handles
  samples as small as 12 observations with structural breaks.
\item
  \textbf{Panel cointegration with breaks and common factors}: Tests by
  Banerjee and Carrion-i-Silvestre (2015) and Banerjee and
  Carrion-i-Silvestre (2025) accommodate structural breaks,
  cross-sectional dependence via common factors, and heterogeneous break
  locations.
\item
  \textbf{Structural break cointegration}: The Hatemi-J (2008) and Maki
  (2012) tests allow for multiple unknown regime shifts in the
  cointegrating relationship.
\item
  \textbf{Fourier-based cointegration}: Tests by Banerjee et al. (2017)
  and Tsong et al. (2016) use trigonometric terms to capture smooth
  structural changes without specifying break dates.
\item
  \textbf{Durbin-Hausman panel tests}: Westerlund (2008) provides tests
  robust to cross-sectional dependence through common factor extraction.
\item
  \textbf{Advanced causality testing}: Beyond standard Granger tests,
  Fourier (Enders and Jones 2016; Nazlioglu et al. 2019), quantile (Cai
  et al. 2023), and bootstrap Fourier quantile (Cheng et al. 2021)
  causality tests offer richer causal inference.
\end{enumerate}

\section{Packages}\label{packages}

\subsection{cointsmall}\label{cointsmall}

Implements the Trinh (2022) cointegration test for very small samples
(as few as 12 observations). Applies residual-based ADF-type tests with
up to two endogenous structural breaks in the constant or both constant
and slope. Includes automatic lag selection via BIC and simulation-based
critical values.

\begin{Shaded}
\begin{Highlighting}[]
\FunctionTok{library}\NormalTok{(cointsmall)}
\NormalTok{result }\OtherTok{\textless{}{-}} \FunctionTok{cointsmall}\NormalTok{(y, x, }\AttributeTok{model =} \StringTok{"cs"}\NormalTok{, }\AttributeTok{max\_breaks =} \DecValTok{2}\NormalTok{)}
\FunctionTok{summary}\NormalTok{(result)}
\end{Highlighting}
\end{Shaded}

\subsection{xtcadfcoint}\label{xtcadfcoint}

Implements the Banerjee and Carrion-i-Silvestre (2025) panel
cointegration test with structural instabilities. Tests the null of no
cointegration using cross-sectionally augmented Dickey-Fuller (CADF)
regressions with Common Correlated Effects (CCE) estimation, allowing
for structural breaks in cointegrating vectors, factor loadings, and
deterministic components.

\begin{Shaded}
\begin{Highlighting}[]
\FunctionTok{library}\NormalTok{(xtcadfcoint)}
\NormalTok{result }\OtherTok{\textless{}{-}} \FunctionTok{xtcadfcoint}\NormalTok{(y }\SpecialCharTok{\textasciitilde{}}\NormalTok{ x1 }\SpecialCharTok{+}\NormalTok{ x2, }\AttributeTok{data =}\NormalTok{ panel\_data,}
                      \AttributeTok{id =} \StringTok{"country"}\NormalTok{, }\AttributeTok{time =} \StringTok{"year"}\NormalTok{, }\AttributeTok{model =} \StringTok{"break"}\NormalTok{)}
\FunctionTok{summary}\NormalTok{(result)}
\end{Highlighting}
\end{Shaded}

\subsection{xtbreakcoint}\label{xtbreakcoint}

Implements panel cointegration tests with structural breaks and
cross-section dependence following Banerjee and Carrion-i-Silvestre
(2015). Provides iterative factor-break estimation, individual ADF tests
on defactored residuals, standardized panel statistics, and the Bai and
Ng (2004) MQ test for identifying common stochastic trends. Five model
specifications with varying deterministic components.

\begin{Shaded}
\begin{Highlighting}[]
\FunctionTok{library}\NormalTok{(xtbreakcoint)}
\NormalTok{result }\OtherTok{\textless{}{-}} \FunctionTok{xtbreakcoint}\NormalTok{(y }\SpecialCharTok{\textasciitilde{}}\NormalTok{ x1 }\SpecialCharTok{+}\NormalTok{ x2, }\AttributeTok{data =}\NormalTok{ panel\_data,}
                       \AttributeTok{id =} \StringTok{"country"}\NormalTok{, }\AttributeTok{time =} \StringTok{"year"}\NormalTok{,}
                       \AttributeTok{model =} \DecValTok{3}\NormalTok{, }\AttributeTok{max\_breaks =} \DecValTok{2}\NormalTok{)}
\FunctionTok{summary}\NormalTok{(result)}
\end{Highlighting}
\end{Shaded}

\subsection{hatemicoint}\label{hatemicoint}

Implements the Hatemi-J (2008) cointegration test with two unknown
regime shifts. Provides ADF\emph{, Zt}, and Zα* test statistics with
endogenously determined break dates. Critical values are based on the
simulation tables of Hatemi-J (2008).

\begin{Shaded}
\begin{Highlighting}[]
\FunctionTok{library}\NormalTok{(hatemicoint)}
\NormalTok{result }\OtherTok{\textless{}{-}} \FunctionTok{hatemicoint}\NormalTok{(y, x, }\AttributeTok{model =} \StringTok{"cs"}\NormalTok{)}
\FunctionTok{summary}\NormalTok{(result)}
\end{Highlighting}
\end{Shaded}

\subsection{xtdhcoint}\label{xtdhcoint}

Implements the Westerlund (2008) Durbin-Hausman panel cointegration
tests, which are robust to cross-sectional dependence through common
factor extraction via principal components. Provides both group-mean
(DHg) and panel (DHp) test statistics with automatic factor number
selection.

\begin{Shaded}
\begin{Highlighting}[]
\FunctionTok{library}\NormalTok{(xtdhcoint)}
\NormalTok{result }\OtherTok{\textless{}{-}} \FunctionTok{xtdhcoint}\NormalTok{(y }\SpecialCharTok{\textasciitilde{}}\NormalTok{ x1 }\SpecialCharTok{+}\NormalTok{ x2, }\AttributeTok{data =}\NormalTok{ panel\_data,}
                    \AttributeTok{id =} \StringTok{"country"}\NormalTok{, }\AttributeTok{time =} \StringTok{"year"}\NormalTok{)}
\FunctionTok{summary}\NormalTok{(result)}
\end{Highlighting}
\end{Shaded}

\subsection{fcoint}\label{fcoint}

Implements four Fourier-based cointegration tests for smooth structural
breaks: the Fourier ADL test (FADL) of Banerjee et al. (2017), the
Fourier Engle-Granger tests (FEG and FEG2), and the Tsong et al. (2016)
KPSS-type test. Lag and frequency selection via AIC or BIC.

\begin{Shaded}
\begin{Highlighting}[]
\FunctionTok{library}\NormalTok{(fcoint)}
\NormalTok{result }\OtherTok{\textless{}{-}} \FunctionTok{fcoint}\NormalTok{(y, x, }\AttributeTok{test =} \StringTok{"fadl"}\NormalTok{, }\AttributeTok{max\_freq =} \DecValTok{3}\NormalTok{, }\AttributeTok{ic =} \StringTok{"aic"}\NormalTok{)}
\FunctionTok{summary}\NormalTok{(result)}
\end{Highlighting}
\end{Shaded}

\subsection{makicoint}\label{makicoint}

Implements the Maki (2012) cointegration test allowing for an unknown
number of structural breaks (up to five). Extends the Gregory-Hansen
framework to multiple breaks with four model specifications: level
shift, level shift with trend, regime shift, and regime shift with
trend.

\begin{Shaded}
\begin{Highlighting}[]
\FunctionTok{library}\NormalTok{(makicoint)}
\NormalTok{result }\OtherTok{\textless{}{-}} \FunctionTok{makicoint}\NormalTok{(y, x, }\AttributeTok{model =} \StringTok{"regime\_shift"}\NormalTok{,}
                    \AttributeTok{max\_breaks =} \DecValTok{3}\NormalTok{, }\AttributeTok{trim =} \FloatTok{0.15}\NormalTok{)}
\FunctionTok{summary}\NormalTok{(result)}
\end{Highlighting}
\end{Shaded}

\subsection{caustests}\label{caustests}

A comprehensive suite of Granger causality tests: standard Toda-Yamamoto
(Toda and Yamamoto 1995), Fourier single-frequency (Enders and Jones
2016) and cumulative-frequency (Nazlioglu et al. 2019) tests, quantile
causality (Cai et al. 2023), and bootstrap Fourier Granger causality in
quantiles (Cheng et al. 2021). All tests include bootstrap inference for
robust p-values.

\begin{Shaded}
\begin{Highlighting}[]
\FunctionTok{library}\NormalTok{(caustests)}
\CommentTok{\# Fourier Toda{-}Yamamoto test}
\NormalTok{result }\OtherTok{\textless{}{-}} \FunctionTok{ftycaus}\NormalTok{(y, x, }\AttributeTok{max\_lag =} \DecValTok{4}\NormalTok{, }\AttributeTok{max\_freq =} \DecValTok{3}\NormalTok{, }\AttributeTok{nboot =} \DecValTok{1000}\NormalTok{)}
\FunctionTok{summary}\NormalTok{(result)}

\CommentTok{\# Quantile causality test}
\NormalTok{result }\OtherTok{\textless{}{-}} \FunctionTok{qcaus}\NormalTok{(y, x, }\AttributeTok{taus =} \FunctionTok{c}\NormalTok{(}\FloatTok{0.25}\NormalTok{, }\FloatTok{0.5}\NormalTok{, }\FloatTok{0.75}\NormalTok{), }\AttributeTok{nboot =} \DecValTok{500}\NormalTok{)}
\FunctionTok{summary}\NormalTok{(result)}
\end{Highlighting}
\end{Shaded}

\section{Acknowledgements}\label{acknowledgements}

The author acknowledges Anindya Banerjee and Josep Lluís
Carrion-i-Silvestre for making their original GAUSS code available, and
Joakim Westerlund for the original GAUSS implementation of the
Durbin-Hausman test.

\section*{References}\label{references}
\addcontentsline{toc}{section}{References}

\protect\phantomsection\label{refs}
\begin{CSLReferences}{1}{1}
\bibitem[\citeproctext]{ref-BaiNg2004}
Bai, Jushan, and Serena Ng. 2004. {``A {PANIC} Attack on Unit Roots and
Cointegration.''} \emph{Econometrica} 72 (4): 1127--77.
\url{https://doi.org/10.1111/j.1468-0262.2004.00528.x}.

\bibitem[\citeproctext]{ref-Banerjee2015}
Banerjee, Anindya, and Josep Lluís Carrion-i-Silvestre. 2015.
{``Cointegration in Panel Data with Structural Breaks and Cross-Section
Dependence.''} \emph{Journal of Applied Econometrics} 30 (1): 1--22.
\url{https://doi.org/10.1002/jae.2348}.

\bibitem[\citeproctext]{ref-Banerjee2024}
Banerjee, Anindya, and Josep Lluís Carrion-i-Silvestre. 2025. {``Panel
Cointegration Tests with Structural Instabilities.''} \emph{Journal of
Business \& Economic Statistics} 42 (3): 1043--55.
\url{https://doi.org/10.1080/07350015.2024.2352073}.

\bibitem[\citeproctext]{ref-Banerjee2017fadl}
Banerjee, Parantap, Vladimir Arcabic, and Hyejin Lee. 2017. {``Testing
for Cointegration in a {Fourier} Autoregressive Distributed Lag
Framework.''} \emph{Economic Modelling} 61: 120--31.
\url{https://doi.org/10.1016/j.econmod.2017.03.004}.

\bibitem[\citeproctext]{ref-Cai2023}
Cai, Yifei, Dayong Zhang, Tsangyao Chang, and Chia-Hao Lee. 2023.
{``Quantile Dependence Between Crude Oil and {China}'s Biofuel Feedstock
Commodity Market.''} \emph{Finance Research Letters} 54: 104327.
\url{https://doi.org/10.1016/j.frl.2023.104327}.

\bibitem[\citeproctext]{ref-Cheng2021}
Cheng, Shu-Yu, Zhe Cao, and Tie Li. 2021. {``Bootstrap {Fourier} Granger
Causality Test in Quantiles.''} \emph{Journal of Statistical Theory and
Practice} 15: 63. \url{https://doi.org/10.1007/s12076-020-00263-0}.

\bibitem[\citeproctext]{ref-Enders2016}
Enders, Walter, and Paul Jones. 2016. {``A Unit Root Test Using a
{Fourier} Series to Approximate Smooth Breaks.''} \emph{Studies in
Nonlinear Dynamics \& Econometrics} 20 (3): 281--301.
\url{https://doi.org/10.1515/snde-2014-0101}.

\bibitem[\citeproctext]{ref-HatemiJ2008}
Hatemi-J, Abdulnasser. 2008. {``Tests for Cointegration with Two Unknown
Regime Shifts with an Application to Financial Market Integration.''}
\emph{Empirical Economics} 35 (3): 497--505.
\url{https://doi.org/10.1007/s00181-007-0175-9}.

\bibitem[\citeproctext]{ref-Maki2012}
Maki, Daiki. 2012. {``Tests for Cointegration Allowing for an Unknown
Number of Breaks.''} \emph{Economic Modelling} 29 (5): 2011--15.
\url{https://doi.org/10.1016/j.econmod.2012.04.022}.

\bibitem[\citeproctext]{ref-Nazlioglu2019}
Nazlioglu, Saban, Ugur Soytas, and Rangan Gupta. 2019. {``Oil Prices and
Financial Stress: A Volatility Spillover Analysis.''} \emph{Emerging
Markets Finance and Trade} 55 (12): 2809--26.
\url{https://doi.org/10.1080/1540496X.2018.1434072}.

\bibitem[\citeproctext]{ref-Pfaff2008}
Pfaff, Bernhard. 2008. \emph{Analysis of Integrated and Cointegrated
Time Series with r}. 2nd ed. Springer.
\url{https://doi.org/10.1007/978-0-387-75967-8}.

\bibitem[\citeproctext]{ref-TodaYamamoto1995}
Toda, Hiro Y., and Taku Yamamoto. 1995. {``Statistical Inference in
Vector Autoregressions with Possibly Integrated Processes.''}
\emph{Journal of Econometrics} 66 (1--2): 225--50.
\url{https://doi.org/10.1016/0304-4076(94)01616-8}.

\bibitem[\citeproctext]{ref-Trinh2022}
Trinh, Thi Huyen. 2022. {``Cointegration Tests for Very Small
Samples.''} \url{https://www.thema.u-cergy.fr/IMG/pdf/2022-01.pdf}.

\bibitem[\citeproctext]{ref-Tsong2016}
Tsong, Ching-Chuan, Cheng-Feng Lee, Li-Ju Tsai, and Te-Chung Hu. 2016.
{``The {Fourier} Approximation and Testing for the Null of
Cointegration.''} \emph{Empirical Economics} 51 (3): 1085--113.
\url{https://doi.org/10.1007/s00181-015-0921-5}.

\bibitem[\citeproctext]{ref-Westerlund2008}
Westerlund, Joakim. 2008. {``Panel Cointegration Tests of the {Fisher}
Hypothesis.''} \emph{Journal of Applied Econometrics} 23 (2): 193--233.
\url{https://doi.org/10.1002/jae.963}.

\end{CSLReferences}
